% Editable workshop manuscript. Source recovered in LA-001; related work
% integrated in LA-021; results, limitations, and conclusions completed in
% LA-022 from admitted automated evidence. See ../source_recovery.md.
\documentclass{article}
\PassOptionsToPackage{numbers,compress}{natbib}
\makeatletter
\def\input@path{{../../../}}
\makeatother
\usepackage{neurips_2026_vericode}
\usepackage[utf8]{inputenc}
\usepackage[T1]{fontenc}
\usepackage{hyperref}
\usepackage{url}
\usepackage{booktabs}
\usepackage{amsfonts}
\usepackage{nicefrac}
\usepackage{microtype}
\usepackage{xcolor}
\usepackage{longtable,array}
\hypersetup{pdftitle={From Law to Action: Neuro-Symbolic Runtime Enforcement for MCP Agents},pdfauthor={}}
\title{From Law to Action: Neuro-Symbolic Runtime Enforcement for MCP Agents}
\workshoptitle{NeurIPS 2026 Workshop on AI for Verifiable Coding}
\author{Anonymous Author(s)}
\begin{document}
\maketitle
\begin{abstract}
We evaluate a bounded runtime-checking question for tool effects that stand in for generated-code side effects: under explicit \texttt{ENFORCE} at one pre-invocation Model Context Protocol (MCP) handler, can an independently observed sandbox mutation that lacks a matching context-bound admission be prevented, while a matched admitted case still performs useful permitted work? Legal documents, CVE-linked vulnerability and repair records, and SkillCenter procedures compile into Legal, Security, and Intent IR. Those IRs are models with distinct authority, not proofs of legal correctness, universal vulnerability absence, or permission. On an admitted $60\times 5\times 3$ matrix of frozen modeled-policy probes, full enforcement observed $0/90$ policy-forbidden handler effects and completed $90/90$ allowed useful work. Unguarded execution and model-free prompt/retrieval labels observed $90/90$ forbidden effects. Lightweight policy plus UCAN leaked $33/90$. Expert legal fidelity, human agreement, and closed-loop model planning were not collected or executed and are withdrawn rather than filled. Closest systems---reference monitors, proof-carrying authorization, Catala, UCAN/Cedar, AgentSpec, Progent, CaMeL, and AgentDojo---motivate implementable arm classes; none is a scored baseline here.
\end{abstract}

\section{Introduction}

An autonomous agent can follow a useful procedure and still violate a legal restriction, reproduce a known vulnerability, or act outside its delegated authority. A safeguard placed only in the model prompt cannot distinguish these cases reliably. The implementation described here places source-grounded constraints and a separate authorization boundary between proposed actions and a protected MCP handler. MCP supplies tool discovery and invocation; application implementations remain responsible for enforcing security requirements \citep{mcp,mcp-auth}.

The method begins before a tool call. A legal collection supplies normative constraints, a vulnerability-and-fix collection supplies security evidence, and a skill collection supplies descriptions of intended procedures. These become three intermediate representations (IRs), not one undifferentiated retrieval index. Formalization produces typed logic views and proof obligations. The runtime joins the results against one exact invocation, then checks authorization again before the effect. This design combines semantic compilation, proof search, capability security, and stateful execution rather than asking a language model to certify its own behavior.

The evaluated workshop contribution is that bounded mediation test \citep{vericode2026cfp}. It is not a new authorization logic, a new legal language, or a first agent runtime monitor. AgentSpec, Progent, and CaMeL already supply runtime rules, tool-privilege checks, and capability isolation \citep{wang2025agentspec,shi2025progent,debenedetti2025camel}; Catala compiles expert-authored law \citep{catala}; AgentDojo measures utility beside attacks \citep{debenedetti2024agentdojo}. We cite those primary sources as comparison classes. We do not execute them as experimental arms. Source adapters, constraint compilation, DuckDB-backed supervision, and content-addressed evidence are implementation context. Legal IR is a lossy model of selected facets, not a proof of legal correctness. Default \texttt{OFF} pass-through and in-memory consumption are not enforcement. The scored study uses one sandbox handler, independently observed filesystem-journal effects, and a frozen modeled policy. It does not evaluate model-generated programs, expert legal interpretation, or universal prevention.

\section{Related work}
\label{sec:related}

The scored evaluation asks whether an independently observed handler effect is denied at one explicit \texttt{ENFORCE} delegate while matched admitted work remains executable. The comparison below is a literature comparison on that question. No cited external system was executed as a scored arm.

\paragraph{Runtime enforcement.}
Complete mediation requires a check at every access to a protected resource \citep{saltzer1975protection}. Security automata and inlined or editing monitors describe which policies a runtime mechanism can enforce \citep{schneider2000enforceable,erlingsson2000irm,ligatti2005edit,hamlen2006enforcement}. Those mechanisms inspect program actions. They do not compile legal or vulnerability sources, and they do not count useful MCP-handler work. AgentSpec supplies agent runtime triggers, predicates, and enforcement rules \citep{wang2025agentspec}. Progent checks symbolic tool privileges and classifies policy updates using an SMT solver \citep{shi2025progent}. These direct precedents rule out claims that this paper introduces agent runtime rules or guarded privilege updates. Neither system was benchmarked here.

\paragraph{Proof-carrying admission.}
Proof-carrying code \citep{necula1997pcc} and proof-carrying authentication \citep{appel1999pca} require a checker to accept a proof before running untrusted code or granting a request. Translation validation and CompCert separate a translator from a checked proof \citep{pnueli1998translation,leroy2009compcert}. A content identifier or a \texttt{proof\_cid} in an envelope is not that check.

\paragraph{Executable law is not checked derivation.}
Encoding a statute as a logic program already performs an interpretation \citep{sergot1986bna,prakken2015lawlogic,benzmueller2020logikey}. Catala compiles literate legal specifications to programs and has been applied to tax law \citep{catala,merigoux2021tax}. It is the closest executable-law precedent. It is not an agent runtime. Legal IR in this paper is a bounded, lossy compilation of selected facets; unsupported facets and expert applicability remain labeling problems, not theorems.

\paragraph{Capabilities and application-level MCP enforcement.}
Object capabilities, attenuating caveats, relation tuples, and analyzable policy languages \citep{miller2006robust,birgisson2014macaroons,pang2019zanzibar,cutler2024cedar} justify a lightweight grant check. UCAN specifies signed delegated capabilities \citep{ucan}. MCP supplies tool invocation; its authorization document specifies OAuth~2.1 resource-server roles and requires implementations to follow OAuth~2.1 security best practices \citep{mcp,mcp-auth}. That protocol is not a generated-code effect monitor. A signature, a CID, or a peer identity is not a generated-code effect observation.

\paragraph{Verifiable coding and agent-security evaluations.}
Functional correctness and issue resolution \citep{chen2021codex,jimenez2024swebench} measure useful coding work without gated forbidden-effect counters. Closed-loop verified generation and autoformalization \citep{sun2024clover,yang2023leandojo,wu2022autoformalization} target specs and prover artifacts, not this handler route. AgentDojo, ToolEmu, InjecAgent, $\tau$-bench, and AgentHarm \citep{debenedetti2024agentdojo,ruan2024toolemu,zhan2024injecagent,yao2024taubench,andriushchenko2025agentharm} measure tool-use attacks, emulated risk, injection, policy following, or harm. CaMeL isolates agent data flow with capabilities \citep{debenedetti2025camel}. They motivate reporting utility beside forbidden effects. They were not run here.

\paragraph{Implementable arms, not imported suites.}
The matched arms are mechanism classes in this stack: an unguarded handler (A0); inert prompt and retrieval labels as model-free equivalence controls (A1, A2); a qualified policy and capability verifier (A3); and explicit \texttt{ENFORCE} with exact context binding, a qualified proof route, and durable consumption (A4). Prompt or retrieval efficacy is not identified in the model-free study. CVEfixes is a vulnerability/fix collection \citep{cvefixes}, not a runtime. DuckDB/Quack, IPFS, and libp2p \citep{quack,ipfs,libp2p} are operational substrates. The workshop fit is runtime checking of a proposed effect before a protected mutation \citep{vericode2026cfp}, not a claim that every cited system was benchmarked.

\section{Corpus sources and compilation into IRs}

The evaluated source freeze is 30 lineage families and 60 cases: six legal documents, 12 CVE repository families with vulnerable/fixed pairs, and 12 SkillCenter families with a base procedure and an adversarial variant. Physical reads used to reserve that freeze were six legal documents, 4329 rows from the first CVEfixes shard, and 1182 SkillCenter index/content rows in the read bundle. Release-card totals (municipal harvests, national-legislation cards, 12{,}987 CVEfixes input rows, 216{,}972 SkillCenter physical records) remain unrecounted background and are not evaluated populations. Appendix~A records family-level provenance requirements. No worldwide-coverage claim follows from the bounded sample.

Legal documents compile through a restricted seven-field grammar (modality, actor, action, object, conditions, exceptions, temporal qualifiers) with source maps and declared losses. Compilation is inspectable. Source-to-rule fidelity remains a separate validation requirement and was not expert-reviewed in this study. Executable-law work such as Catala supplies a relevant precedent \citep{catala}; it is not an agent runtime or a handler-effect observer.

CVEfixes links vulnerability records to source repositories and fixes \citep{cvefixes}. The inspected release pins hitoshura25/cvefixes at \path{d4f5c4ea65329d9ccbb8a3b3149e5d06eda5edb2}. The adapter constructs a canonical Security IR declaration and rejects wildcard scope and inserted authority grants. Vulnerable code is a positive control for a candidate restriction; fixed code is a negative control for that restriction. A fix does not certify that all other vulnerabilities are absent. Runtime \texttt{cve\_security\_gate} maps declared intent and independently extracted handler effects into the same Security IR root. Incomplete mappings remain unknown.

SkillCenter supplies procedural Markdown. The bounded reader opens upstream SQLite bundles read-only, disables extension loading, and never executes commands appearing in a skill. The inspected source is Tommysha/skillcenter-bundles at \path{f9dd4fec3c86d85ebf116c7408ac5ce602c418a1}. Intent IR distinguishes goals, guards, effects, and control edges, with source spans and grounded-versus-inferred status. A procedure's word ``permitted'' is a source assertion, not a capability. The default skill path is a deterministic structural normalizer. No trained encoder was pinned or evaluated.

The three families share identity, provenance, and obligation contracts. They do not share an ontology or grant authority to one another. Sparse, dense, and graph retrieval may nominate evidence; hard applicability and exact identity determine whether it may be used. Retrieval engineering is not a measured token-saving result.
\begingroup
\renewcommand{\thetable}{1}
\begin{longtable}{@{}>{\raggedright\arraybackslash}p{.24\linewidth} >{\raggedright\arraybackslash}p{.32\linewidth} >{\raggedright\arraybackslash}p{.38\linewidth}@{}}
\caption{Three source-to-IR paths. None of the source streams directly supplies execution authority.}\label{tab:1}\\
\toprule
\textbf{Source stream} & \textbf{Derived semantics} & \textbf{Contribution to an action decision} \\
\midrule
\endfirsthead
\textbf{Source stream} & \textbf{Derived semantics} & \textbf{Contribution to an action decision} \\
\midrule
\endhead
\bottomrule
\endlastfoot
Legal texts and versions & Norms, scope, authority, exceptions, temporal applicability & Which represented legal constraints apply? \\[3pt]
CVE-linked code and fixes & Threat assumptions, effects, mitigations, scoped prohibition candidates & Does the proposed or generated behavior violate a reviewed security constraint? \\[3pt]
Skill Markdown and metadata & Goals, actions, guards, effects, verification, control flow & What is the agent trying to do, and what effects does its procedure require? \\[3pt]
\end{longtable}
\endgroup

\section{Logic families, joining, and the protected route}

An IR is a domain representation; a logic family determines how its statements are interpreted; a prover is a backend that checks a supported encoding. A family's presence in a registry does not mean every formula or backend is qualified. Appendix~B records the broader map. The benchmark-required proof-oriented route actually executed is a QF\_BOOL / propositional SAT encoding of a Hoare-style security invariant, with SymPy DPLL as a child process and an independent exhaustive truth-table checker. Both emit satisfiability authority only. SAT-only, simulated, policy-only, and monitor evidence cannot satisfy a \texttt{theorem\_proof} requirement. QF\_LIA interpolation and interactive kernels were not selected and were not available on the sealed validation path.
\begingroup
\renewcommand{\thetable}{2}
\begin{longtable}{@{}>{\raggedright\arraybackslash}p{.24\linewidth} >{\raggedright\arraybackslash}p{.32\linewidth} >{\raggedright\arraybackslash}p{.38\linewidth}@{}}
\caption{Roles of native views. Only the QF\_BOOL SAT route was qualified for the scored study.}\label{tab:2}\\
\toprule
\textbf{Logic or formal view} & \textbf{Implemented role in the stack} & \textbf{Qualified in this evaluation} \\
\midrule
\endfirsthead
\textbf{Logic or formal view} & \textbf{Implemented role in the stack} & \textbf{Qualified in this evaluation} \\
\midrule
\endhead
\bottomrule
\endlastfoot
Hoare-style contracts and SAT & Preconditions, postconditions, and a propositional invariant over grant and forbidden effect & QF\_BOOL SAT plus independent truth table; satisfiability only \\[3pt]
Rule/policy and capability views & Bounded authorization declarations and signed UCAN grants & Real Ed25519 verification in A3/A4; not a theorem \\[3pt]
Deontic / temporal / event-calculus tags & Native representations and compiler bridges & Registered; not a live scored backend \\[3pt]
Interactive proof kernels & Lean/Coq/Isabelle adapters & Unavailable on the sealed path; no invented kernel result \\[3pt]
\end{longtable}
\endgroup

An InvocationIntentEnvelope commits to source identity, tenant, actor, audience, tool/version, arguments, proposed effects, environment, nonce, and relevant roots. Trusted adapters bind the effect description to the actual operation. The constraint compiler joins Legal, Security, and Intent checks against that exact invocation. A favorable score cannot offset a failed hard constraint.

The selected protected route is \path{SupervisorPreInvocationEnforcement.authorize_and_delegate} in explicit \texttt{ENFORCE} around \path{BoundedExportHandler.export_json}. Independent filesystem-journal counters observe handler effects. A permissible candidate still requires a current grant for the exact actor, tool, resource, and bounds. The runtime revalidates and consumes the permitted use immediately before delegation. Replay, expiry, wrong audience, changed arguments, undeclared effects, and changed roots are denials at that boundary. DuckDB with a typed Quack owner records operational progression for the A4 deployment; the default in-memory store is a labeled contrast only. IPFS names immutable bytes \citep{ipfs}; libp2p carries framed MCP messages \citep{libp2p}; UCAN checks signed attenuated delegation \citep{ucan}. None of those substrates is itself the scored safety result. Kit's \texttt{AuthorizationGate} is a related policy gate and is not the selected safety claim.

% Results consumed from LA-029 admitted summary/cost/claim guidance, LA-009
% source-contract metrics, and explicit scoped omissions. No invented cells.
\section{Evaluation}
\label{sec:eval}

\subsection{Frozen protocol and automated evidence scope}
\label{sec:eval-protocol}

The predeclared fixed-action study uses 30 source-lineage families and 60 cases (six legal families, 12 CVE repository pairs, 12 SkillCenter families), split 6/6/18 into development, calibration, and final by source lineage \citep{cvefixes}. Each family contributes one policy-allowed and one policy-forbidden case. Five matched arms share the same candidate identifiers, arguments, and sandbox handler: unguarded execution (A0); inert prompt-text and retrieval-context labels as model-free equivalence controls (A1, A2); lightweight policy plus real Ed25519 UCAN verification (A3); and explicit \texttt{ENFORCE} with exact-context binding, a qualified QF\_BOOL SAT job, and file-backed DuckDB consumption (A4). Seeds $104729$, $104759$, and $104761$ yield a $60\times 5\times 3=900$-cell matrix. The independent unit is the lineage family. Allowed and forbidden labels are the frozen modeled policy and machine-checkable behavior contract, not real-world legality.

Outside reviewers were not a run dependency. Expert legal applicability, exception fidelity, semantic accuracy, inter-annotator agreement, and legal validity were not collected and are not scored. Optional author reading, if it occurs later, is non-independent and is not human gold. Generated interpretations remain predictions. The LA-009 skill-validator comparison of $18/18$ is shared-producer conformance: the same Intent IR decoder and schema functions already invoked by the normalizer, not an independent checker.

\subsection{Source-to-IR contracts}
\label{sec:eval-source}

Every frozen case has a retained prediction, failure, or unsupported record; cohorts were not changed after outcomes. Quoted source-span linkage held on $60/60$ cases. Parse or schema validity held on $50/60$. The scoped machine-contract coverage was $60/60$. Four CVE adapter errors and six skill-policy-blocked unsupported rows remain in the denominator accounting. CVE source-difference observation was $24/24$ under an independent behavior-contract producer; polarity remains unknown on all $24$ pairs, and transfer to the \texttt{export\_json} handler is unsupported. Implementation agreement is not semantic accuracy. These measurements replace the original expert-reviewed legal-fidelity row, which is withdrawn as unmeasured (Table~\ref{tab:E1}).

\subsection{Admitted protected-effect and useful-work results}
\label{sec:eval-admitted}

The primary empirical study is the LA-029 operator matrix: $900$ scientific cells, all resource-qualified, with independent filesystem-journal effect counters. Compact rows do not retain full per-cell proof or capability payloads; pinned executed mechanisms, effect counters, and the shared durable-store identity support the runtime observations. Implementation revision \texttt{ee6d73c4a5fc30d05ec4e787fdeb46d6701a4f947b7c874aecabf603e31ce06c}. Model calls were $0$. Independent human gold is false. Two original false host admissions (schedule indices $459$ and $560$) remain labeled as such. Two container-startup recoveries at indices $591$ and $625$ were charged once after proof that science had not started; scientific retries were $0$. Host attempts total $902$. Measured descendant-inclusive group CPU was $1608.817478$ seconds, under the frozen $18{,}000$-second stop. Peak cell memory was $274812928$ bytes; maximum whole-cell wall time was $18.022665493073873$ seconds.

Table~\ref{tab:3} reports the admitted rates. On the held-out final split ($18$ families; $54$ forbidden-labelled and $54$ allowed-labelled attempts per arm), A4 observed $0/54$ policy-forbidden effects and completed $54/54$ allowed useful work with $0/54$ false denials. Across the full admitted matrix the same pattern holds: A4 $0/90$ forbidden effects and $90/90$ allowed useful work; A0--A2 $90/90$ forbidden effects on forbidden-labelled attempts; A3 $33/90$ forbidden effects. A1 and A2 match A0 (final-split paired family difference $0$, cluster-bootstrap $95\%$ interval $[0,0]$) and therefore remain instrumentation/equivalence controls, not prompt or retrieval efficacy. Final-split paired family differences in forbidden-effect rate are $-0.6666666666666666$ for A3 versus A0 (percentile $95\%$ $[-0.8888888888888888,-0.4444444444444444]$), $-1.0$ for A4 versus A0 ($[-1.0,-1.0]$), and $-0.3333333333333333$ for A4 versus A3 ($[-0.5555555555555556,-0.1111111111111111]$). A degenerate zero interval does not establish population-level absence or universal prevention.

The handler under test accepts a declarative \texttt{export\_json} instruction. Independently observed effects are sandbox filesystem-journal events, not AST extraction from model-generated source. No scientific generative model was pinned.

\begingroup
\renewcommand{\thetable}{3}
\begin{table}[t]
\centering
\caption{Admitted LA-029 fixed-action rates from independent effect counters. Forbidden-effect rate uses forbidden-labelled attempts; allowed-task success uses allowed-labelled attempts. Labels are the frozen modeled policy. Overall $n=90$ forbidden-labelled and $90$ allowed-labelled attempts per arm ($180$ decisions). Final $n=54+54$.}
\label{tab:3}
{\small
\begin{tabular}{@{}lcccc@{}}
\toprule
Arm & Split & Forbidden effects & Allowed useful work & False denials \\
\midrule
A0 unguarded & overall & $90/90$ & $90/90$ & $0/90$ \\
A1 prompt label & overall & $90/90$ & $90/90$ & $0/90$ \\
A2 retrieval label & overall & $90/90$ & $90/90$ & $0/90$ \\
A3 policy+UCAN & overall & $33/90$ & $90/90$ & $0/90$ \\
A4 \texttt{ENFORCE} & overall & $0/90$ & $90/90$ & $0/90$ \\
\midrule
A0 unguarded & final & $54/54$ & $54/54$ & $0/54$ \\
A3 policy+UCAN & final & $18/54$ & $54/54$ & $0/54$ \\
A4 \texttt{ENFORCE} & final & $0/54$ & $54/54$ & $0/54$ \\
\bottomrule
\end{tabular}
}
\end{table}
\endgroup

\subsection{Negative results and scoped omissions}
\label{sec:eval-negative}

A3 is not equivalent to A4. The $33$ A3 leaks (decision \texttt{allow}, observed forbidden effect) concentrate on \texttt{wrong\_date\_or\_jurisdiction} ($9$), \texttt{forged\_receipt} ($6$), \texttt{replay} ($6$), \texttt{changed\_root\_clock\_or\_environment} ($6$), and \texttt{undeclared\_handler\_effect} ($6$). On the final split the residual classes are replay, changed root/clock/environment, and wrong date/jurisdiction ($6$ each). A4 denied those identities. SAT/UNSAT on the selected QF\_BOOL route has satisfiability authority only; it cannot authorize a \texttt{theorem\_proof} allow. Z3, cvc5, Vampire, E, Lean, Coq/Rocq, Isabelle, and QF\_LIA interpolation were absent from the sealed validation path and are not assigned invented results.

Closed-loop generated-code planning (protocol RQ3) is explicitly withdrawn: $900/900$ scheduled cells were \texttt{not\_started}, model calls $0$, generated programs $0$. Unstarted zeros are accounting, not measured agent failure. Results are not pooled with the fixed-action matrix.

The original LA-015 compact log and its rescue remain protocol-unadmitted diagnostics; they are not headline admitted rates. LA-017 paired mechanism ablations and LA-016/LA-017 comparison wording are protocol-unadmitted diagnostic evidence and do not establish admitted model efficacy. LA-018 cold/warm/stale timings are a development A4 instrumentation study ($n=12$ per condition; cold mean end-to-end wall $0.012326$\,s, warm $0.006635$\,s, stale-root and stale-clock $0.00322$\,s with $12/12$ stale-evidence flags; model tokens unobserved; no price conversion) and are not a $900$-cell rescore. LA-020 twelve-step traces are separate diagnostic replays of frozen A4 identities, not additional scored cells. Qualification suites for mediation, DuckDB consumption, and local stdio/HTTP/libp2p parity passed their selected-claim cases and remain qualification, not scored arm rates. Wide-area publication and an IPFS daemon were not available.

Admitted rates are taken from \texttt{results/fixed\_actions\_admitted/summary.json} together with the matching cost report and claim guidance. LA-019 generated tables and figures analyze the historical unadmitted log and are not substituted for that admitted summary.

% Limitations required by the automated evidence amendment and LA-029 claim guidance.
\section{Limitations and modeling assumptions}
\label{sec:limits}

Legal IR is a lossy model of selected facets (modality, actor, action, object, conditions, exceptions, temporal qualifiers), not the law and not a proof of legal correctness. Unsupported constructs must not disappear into a successful admission. Expert legal fidelity, human agreement, and legal-validity rates were not collected; they remain unmeasured rather than filled with compiler agreement or generated labels. Security IR is a scoped candidate prohibition. A CVE or CWE identifier, a fix, or a source-difference observation is not universal vulnerability absence and is not transfer to unrelated code. Skill Markdown is untrusted data; a procedure's word ``permitted'' is not a capability.

Complete mediation is claimed only for \path{SupervisorPreInvocationEnforcement.authorize_and_delegate} in explicit \texttt{ENFORCE} around one benchmark-owned \path{BoundedExportHandler.export_json} sandbox mutation. Default \texttt{OFF} is pass-through. The default in-memory consumption store is not the restart-safety configuration; A4 used a file-backed DuckDB typed Quack owner. A DuckDB commit is not remote exactly-once evidence. Unprotected, observational, and nonidentical routes were exercised in qualification and excluded from the safety claim. Local stdio, loopback HTTP, and loopback libp2p parity do not establish wide-area transport.

The selected proof route is propositional SAT with an independent truth-table checker. Satisfiability is not a kernel theorem. Compact admitted records do not retain full per-cell proof or capability payloads. The evaluated handler consumes a declarative export instruction; closed-loop model-generated programs were not executed. No trained skill normalizer or encoder was evaluated. Release-card corpus totals were not independently recounted and are not evaluated populations. Optional author review is not a prerequisite for this draft, is not independent annotation, and cannot replace the frozen automated evidence.


\section{Conclusion}

On one explicit \texttt{ENFORCE} sandbox route, independently observed forbidden handler effects were prevented for every frozen policy-forbidden probe while matched admitted work still succeeded. Lightweight policy plus UCAN was not sufficient for the residual mutation classes. Unguarded execution and inert prompt or retrieval labels were not sufficient either. The result is a bounded, policy-relative verifiable-coding measurement \citep{vericode2026cfp}, not superiority to Catala, AgentSpec, Progent, CaMeL, AgentDojo, or any other listed system, and not a claim of legal correctness, expert agreement, or closed-loop generated-program control. Those latter studies remain unmeasured or withdrawn. Optional author review is not required to generate this evidence-scoped draft and would not be independent human validation.

\clearpage
% Selected bibliography is related_work.bib. Do not also load references.bib
% (duplicate keys) or \nocite unrun systems as if they were scored arms.
\bibliographystyle{unsrtnat}
\bibliography{related_work}
\clearpage
\appendix

\section{Corpus lineage and release accounting}

The author-maintained legal program includes public releases for the U.S.\ Code, state statutes, state administrative rules, Federal Register documents, state/local court rules, municipal/county codes, caselaw, and national legislation. The paper does not interpret a dataset title, a country registry entry, or a successful scrape as exhaustive coverage. The evaluated legal cohort is the six frozen documents. Card-reported municipal, Netherlands, and Belgian quantities have not been independently recounted and are not evaluated records.
\begingroup
\renewcommand{\thetable}{A1}
\begin{longtable}{@{}>{\raggedright\arraybackslash}p{.24\linewidth} >{\raggedright\arraybackslash}p{.32\linewidth} >{\raggedright\arraybackslash}p{.38\linewidth}@{}}
\caption{Legal source families and provenance requirements. Identifying publisher/dataset mappings belong in the private author record until camera-ready.}\label{tab:A1}\\
\toprule
\textbf{Source family} & \textbf{Representative inspected substrate} & \textbf{Required distinctions} \\
\midrule
\endfirsthead
\textbf{Source family} & \textbf{Representative inspected substrate} & \textbf{Required distinctions} \\
\midrule
\endhead
\bottomrule
\endlastfoot
Federal statutory law & U.S. Code sparse GraphRAG release, pinned source revision and official release point & Codification/release point, provision identity, amendment history and actual effective dates \\[3pt]
State law & State-statute release and separate administrative-rule corpus & State, authority family, rule/statute distinction, missing jurisdictions and time coverage \\[3pt]
Federal regulatory material & Federal Register collection & Proposed rule, final rule, notice, publication date, effective date and codified regulation are not interchangeable \\[3pt]
Court rules and caselaw & Merged per-state court rules; Caselaw Access Project-derived collection & Court, jurisdiction, local/general rule, opinion treatment, holdings versus quotations \\[3pt]
Municipal and county law & Multi-drop municipal harvest & Publisher, jurisdiction, per-drop discovery and failure counts; registry aliases may lag new release namespaces \\[3pt]
International legislation & Netherlands corpus; Belgian official-source archive harvest; Argentine CID-keyed retrieval release & Language, source authority, consolidation date, archive capture, extraction version and incomplete coverage \\[3pt]
\end{longtable}
\endgroup

The security source pins hitoshura25/cvefixes at \path{d4f5c4ea65329d9ccbb8a3b3149e5d06eda5edb2}. Graph nodes, edges, vectors, source rows, and unique CVEs are different populations. The evaluated CVE cohort is 12 repository families / 24 cases plus 84 excluded controls that never enter empirical sample counts. Evaluation includes the original vulnerable condition and the fixed negative control. Transfer to a proposed tool needs an explicit mapping; polarity remains unknown where the machine contract does not support it.

The SkillCenter source is Tommysha/skillcenter-bundles at \path{f9dd4fec3c86d85ebf116c7408ac5ce602c418a1}. SQLite is an untrusted input container, not the supervisor database. The evaluated skill cohort is 12 families / 24 cases. Index totals are not additional skills. Counts of available records do not establish counts of normalized Intent IR documents or safe procedures.

\begin{verbatim}
pinned source revision + raw artifact bytes
  -> bounded source record + content identity + use policy
  -> LegalIR | SecurityIR | IntentIR declaration
  -> native formal views + source map + loss/coverage diagnostics
  -> selected applicable constraints + exact invocation
  -> proof jobs + checked evidence + decision receipt
  -> authenticated capability / execution permit
  -> guarded effect + operational observation + durable record
\end{verbatim}

\section{Logic-family and backend method map}

A shared symbol table does not prove that two encodings preserve meaning. The TDFOL and DCEC Python cores establish parsable native representations. Resource-scoped authorization may be representable without invoking every family. Solver success, kernel acceptance, runtime monitors, and policy approval are non-interchangeable authority kinds. A timeout or unavailable backend is not a refutation.
\begingroup
\renewcommand{\thetable}{B1}
\begin{longtable}{@{}>{\raggedright\arraybackslash}p{.24\linewidth} >{\raggedright\arraybackslash}p{.32\linewidth} >{\raggedright\arraybackslash}p{.38\linewidth}@{}}
\caption{Logic families have different meanings. A shared symbol table or cross-view link does not prove that two encodings preserve those meanings.}\label{tab:B1}\\
\toprule
\textbf{Family/view} & \textbf{Representation and operations} & \textbf{What the paper may claim} \\
\midrule
\endfirsthead
\textbf{Family/view} & \textbf{Representation and operations} & \textbf{What the paper may claim} \\
\midrule
\endhead
\bottomrule
\endlastfoot
Propositional and first-order & Boolean operators, typed terms, predicates, variables and quantifiers & Bounded facts and implications under the encoded vocabulary; not general natural-language understanding \\[3pt]
F-logic / frame & Object identity, class membership, slots and relational inheritance & Structured semantic relations; not a cryptographic or temporal proof by itself \\[3pt]
Deontic & Obligation, permission, prohibition & Explicit normative force under the selected axioms; an encoded permission is not an issued capability \\[3pt]
Temporal / TDFOL & Always, eventually, next, until, since & Model time-sensitive requirements; a backend must support the actual fragment \\[3pt]
Event calculus & Events, fluents, initiation/termination & Reason about transitions as represented; an observed event record is not a proof of every transition \\[3pt]
Datalog / rule-policy & Bounded rules and authorization declarations & Rule derivation with an explicit open/closed-world convention, not automatic theorem authority \\[3pt]
Hoare-style / SAT & Preconditions, postconditions, propositional invariants & The qualified QF\_BOOL SAT route; SAT/UNSAT cannot authorize theorem\_proof allows \\[3pt]
Interactive proof & Native proof terms/scripts & Unavailable here; no invented kernel result \\[3pt]
\end{longtable}
\endgroup

\section{Database, artifact, network, and capability boundaries}
\begingroup
\renewcommand{\thetable}{C1}
\begin{longtable}{@{}>{\raggedright\arraybackslash}p{.24\linewidth} >{\raggedright\arraybackslash}p{.32\linewidth} >{\raggedright\arraybackslash}p{.38\linewidth}@{}}
\caption{Database and protocol responsibilities. Legacy or auxiliary indexes can exist without changing the canonical DuckDB control-plane description.}\label{tab:C1}\\
\toprule
\textbf{Layer} & \textbf{Technology and inspected role} & \textbf{Incorrect inference to avoid} \\
\midrule
\endfirsthead
\textbf{Layer} & \textbf{Technology and inspected role} & \textbf{Incorrect inference to avoid} \\
\midrule
\endhead
\bottomrule
\endlastfoot
Supervisor operational authority & DuckDB control.duckdb; typed Quack owner/client; StateTransaction CAS and idempotency & ``The supervisor uses SQLite'' or ``each agent writes the database file directly'' \\[3pt]
Source input packaging & Read-only SkillCenter SQLite bundles; source Parquet/JSON-LD/HTML/WARC & ``An input database format defines the runtime architecture'' \\[3pt]
Retrieval products & BM25, vectors, graphs and bounded Parquet shards keyed by source IDs/CIDs & ``A retrieved neighbor is a grant or a proof'' \\[3pt]
Immutable artifacts & Kit-managed content-addressed blocks, manifests, maps, proof/decision receipts, optional IPFS replication & ``The newest block is automatically the accepted state'' \\[3pt]
Peer carriage & libp2p negotiated streams and MCP++ framing & ``A peer identity grants tool permission'' \\[3pt]
Delegated capability & Strict signed UCAN verification and revocation/replay state & ``A valid token proves the action is legally and semantically safe'' \\[3pt]
\end{longtable}
\endgroup

A caller supplies a named command and typed parameters, not a SQL string. The identical command can retrieve a committed result after a lost reply; a changed payload with the same identity is a conflict. The Quack gateway extension was unavailable under the sealed home and is not claimed. IPFS answers which artifact bytes are requested; local CID computation proves neither publication nor public availability. libp2p answers how peers connect. In-process framing is not libp2p network evidence. Strict UCAN checks cryptographic issuer identity, delegation continuity, resource/ability bounds, temporal validity, attenuation, revocation, and replay policy. It does not infer delegation from natural-language text.

\section{Implementation trace}

The private companion contains exact repository revisions and paths. Neutral aliases should replace identifying names in a blind artifact.
\begingroup
\renewcommand{\thetable}{D1}
\begin{longtable}{@{}>{\raggedright\arraybackslash}p{.24\linewidth} >{\raggedright\arraybackslash}p{.32\linewidth} >{\raggedright\arraybackslash}p{.38\linewidth}@{}}
\caption{Implementation-oriented trace. Compact LA-029 rows remain the scored outcomes; LA-020 retains intermediate artifacts on a separate diagnostic replay of frozen A4 identities.}\label{tab:D1}\\
\toprule
\textbf{Method} & \textbf{Representative implementation} & \textbf{Inputs inspected in a reproducible run} \\
\midrule
\endfirsthead
\textbf{Method} & \textbf{Representative implementation} & \textbf{Inputs inspected in a reproducible run} \\
\midrule
\endhead
\bottomrule
\endlastfoot
Legal routing & Deterministic deontic parser; typed canonical compiler & Dataset revision, source body, jurisdiction, source maps, unsupported facets \\[3pt]
Skill import & SkillCenterSkillRecord and read-only bundle adapter & Physical source row, metadata, skill/library Markdown, entry/body identities \\[3pt]
CVE adaptation & adapt\_\allowbreak{}cvefixes\_\allowbreak{}candidate & Vulnerable/fixed source records, candidate scope, prohibition/control polarity \\[3pt]
Constraint join & IRConstraintCompiler & Candidate graph, legal/security/intent checks, proof requirements, current roots \\[3pt]
Proof authorization & IntentAuthorizationService; SAT child process plus truth-table checker & QF\_BOOL jobs, satisfiability authority, rejected theorem\_proof upgrades \\[3pt]
Guarded execution & Pre-invocation enforcement wrapper in ENFORCE & Current live context, authenticated receipt, expiry/use state, delegate effects \\[3pt]
State control & QuackStateClient; file-backed DuckDB consumption & Named command, generation, fence, revision, idempotency and committed result \\[3pt]
\end{longtable}
\endgroup

A diagnostic replay of the permitted A4 identity recorded by LA-020 (seed 104729, skill case-0; exact attempt identifier retained in the worked-trace record) ran source-pin resolution, Intent IR derivation, exact tool/argument binding, QF\_BOOL SAT (provider unsat, checker agrees, satisfiability only), real Ed25519 UCAN allow, ENFORCE consume on the DuckDB owner, one \texttt{export\_json} dispatch, and one independent journal event with useful work. Matched forbidden replays for wrong audience, undeclared effect, and replay denied at delegation or consumption with zero forbidden effects. LegalIRCompilerAPI was not the selected adapter. Independent human review was not used. No fabricated signature or conceptual step is presented as executed.

\section{Resolved measurement inventory}

Original Table~E1 listed placeholders. Table~\ref{tab:E1} records the actual disposition. ``Withdrawn'' and ``qualification only'' are scoped omissions, not zeros invented from missing runs.
\begingroup
\renewcommand{\thetable}{E1}
\begin{longtable}{@{}>{\raggedright\arraybackslash}p{.24\linewidth} >{\raggedright\arraybackslash}p{.32\linewidth} >{\raggedright\arraybackslash}p{.38\linewidth}@{}}
\caption{Resolved measurement inventory. Withdrawn claims are unmeasured. Qualification suites are not scored arm rates. The primary protected-effect study is the admitted LA-029 matrix.}\label{tab:E1}\\
\toprule
\textbf{Measurement} & \textbf{Actual record} & \textbf{Disposition} \\
\midrule
\endfirsthead
\textbf{Measurement} & \textbf{Actual record} & \textbf{Disposition} \\
\midrule
\endhead
\bottomrule
\endlastfoot
Legal source-to-IR fidelity & Expert-reviewed atoms, exceptions, applicability and agreement & Withdrawn / unmeasured; replaced by automated span, schema, and machine-contract checks ($60/60$, $50/60$, $60/60$) \\[3pt]
CVE restriction specificity & Vulnerable/fixed pairs and excluded controls & 12 pairs / 24 cases and 84 excluded controls retained; source-difference $24/24$; polarity unknown $24/24$; no transfer claim \\[3pt]
Skill grounding & Source spans and malicious Markdown negatives & Span linkage on all skill cases; $18/18$ shared-producer schema conformance, not independent agreement; six unsupported policy-blocked rows \\[3pt]
Proof job execution & Per-family provider/checker route & QF\_BOOL SAT qualified; theorem kernels and QF\_LIA interpolation unavailable, not invented \\[3pt]
Guarded effect prevention & Delegate/effect counters under allow, deny, stale and replay & Admitted LA-029: A4 $0/90$ forbidden effects and $90/90$ allowed useful work \\[3pt]
DuckDB owner safety & Wrong generation/fence, concurrent commands, lost reply, restart & Qualification $14/14$ selected-claim cases; in-memory restart is a labeled contrast; not a scored $900$-cell rate \\[3pt]
Cross-transport parity & Stdio/HTTP/P2P same request and outcome & Local-process and loopback qualification; WAN and IPFS daemon absent \\[3pt]
End-to-end efficiency & Model tokens, retrieval, proof and state overhead & Development A4 instrumentation only; model tokens unobserved; no price conversion; not a $900$-cell rescore \\[3pt]
Closed-loop planning & Pinned-model generated programs and recovery & Explicitly withdrawn: $900/900$ not started; $0$ model calls; $0$ generated programs \\[3pt]
\end{longtable}
\endgroup

\section{Scope, ongoing work, and disclosure}

Semantic-addressed program worlds, learned residuals, semantic-preserving refactoring, parallel sealing, and proof-carrying test reuse were unavailable for this revision and are not implemented contributions. Semantic minification and learned encoders were not evaluated as cost savings.

This draft was prepared using language-model assistance for source inspection, organization, and writing. Manuscript generation and the evidence-scoped submission candidate require no outside review. Optional author review is labeled non-independent and is not a prerequisite. No independent human legal, security, or intent validation and no inter-annotator agreement study was collected. Exact tools, model versions, and executed experiments for camera-ready disclosure remain an LA-024 obligation; the scientific claims above are bound to the retained receipts cited in the claim--evidence matrix.

\end{document}
